\section{AIME I 1988}
        \begin{enumerate}
        	\item (\textit{AIME I 1988}) One commercially available ten-button lock may be opened by depressing -- in any order -- the correct five buttons. The sample shown below has $\{1, 2, 3, 6, 9\}$ as its combination. Suppose that these locks are redesigned so that sets of as many as nine buttons or as few as one button could serve as combinations. How many additional combinations would this allow?


 
                    \begin{figure}[H]
                        \centering
                        \includegraphics[width=0.5\textwidth]{Problem/AIME_I_1988_1_Problem_Figure_1.png}
                    \end{figure}
                    

 \textbf{Answer: }770



\textbf{Solution}

Currently there are ${10 \choose 5}$ possible combinations.With any integer $x$ from $1$ to $9$, the number of ways to choose a set of $x$ buttons is $\sum^{9}_{k=1}{10 \choose k}$.Now we can use the identity $\sum^{n}_{k=0}{n \choose k}=2^{n}$.So the number of additional combinations is just$2^{10}-{10\choose 0}-{10\choose 10}-{10 \choose 5}=1024-1-1-252=\boxed{770}$.

 
Note: A simpler way of thinking to get $2^{10}$ is thinking that each button has two choices, to be black or

 to be white.

 


	\item (\textit{AIME I 1988}) For any positive integer $k$, let $f_1(k)$ denote the square of the sum of the digits of $k$.  For $n \ge 2$, let $f_n(k) = f_1(f_{n - 1}(k))$.  Find $f_{1988}(11)$.


 
 \textbf{Answer: }169



\textbf{Solution}

We see that $f_{1}(11)=4$
$f_2(11) = f_1(4)=16$
$f_3(11) = f_1(16)=49$
$f_4(11) = f_1(49)=169$
$f_5(11) = f_1(169)=256$
$f_6(11) = f_1(256)=169$

 Note that this revolves between the two numbers. Since $1988$ is even, we thus have $f_{1988}(11) = f_{4}(11) = \boxed{169}$.

 


	\item (\textit{AIME I 1988}) Find $(\log_2 x)^2$ if $\log_2 (\log_8 x) = \log_8 (\log_2 x)$.


 
 \textbf{Answer: }027



\textbf{Solution 1}

Raise both as \href{https://artofproblemsolving.com/wiki/index.php?title=Exponent}{exponents} with base 8:

 
 \begin{align*} 8^{\log_2 (\log_8 x)} &= 8^{\log_8 (\log_2 x)}\\ 2^{3 \log_2(\log_8x)} &= \log_2x\\ (\log_8x)^3 &= \log_2x\\ \left(\frac{\log_2x}{\log_28}\right)^3 &= \log_2x\\ (\log_2x)^2 &= (\log_28)^3 = \boxed{027}\\ \end{align*}


 A quick explanation of the steps: On the 1st step, we use the property of \href{https://artofproblemsolving.com/wiki/index.php?title=Logarithm}{logarithms} that $a^{\log_a x} = x$. On the 2nd step, we use the fact that $k \log_a x = \log_a x^k$. On the 3rd step, we use the \href{https://artofproblemsolving.com/wiki/index.php?title=Change\_of\_base\_formula}{change of base formula}, which states $\log_a b = \frac{\log_k b}{\log_k a}$ for arbitrary $k$.

 



\textbf{Solution 2: Substitution}

We wish to convert this expression into one which has a uniform base. Let's scale down all the powers of 8 to 2. 

 
 \begin{align*} {\log_2 (\frac{1}{3}\log_2 x)} &= \frac{1}{3}{\log_2 (\log_2 x)}\\ {\log_2 x = y}\\ {\log_2 (\frac{1}{3}y)} &= \frac{1}{3}{\log_2 (y)}\\ {3\log_2 (\frac{1}{3}y)} &= {\log_2 (y)}\\  {\log_2 (\frac{1}{3}y)^3} &= {\log_2 (y)}\\  \end{align*}
Solving, we get $y^2 = 27$, which is what we want. $\boxed{27}$

 


 Just a quick note- In this solution, we used 2 important rules of logarithm:1) $\log_a b^n=n\log_a b$.2) $\log_{a^n} b=\frac{1}{n}\log_a b$.

 



\textbf{Solution 3}

First we have
 \begin{align*} \log_2(\log_8x)&=\log_8(\log_2x)\\ \frac{\log_2(\log_8x)}{\log_8(\log_2x)}&=1 \end{align*}
Changing the base in the numerator yields
 \begin{align*} \frac{3\log_8(\log_8x)}{\log_8(\log_2x)}&=1\\ \frac{\log_8(\log_8x)}{\log_8(\log_2x)}&=\frac{1}{3}\\ \end{align*}
Using the property $\frac{\log_ab}{\log_ac}=\log_cb$ yields
 \begin{align*} \log_{\log_2x}(\log_8x)&=\frac{1}{3}\\ (\log_2x)^\frac{1}{3}&=\log_8x\\ \sqrt[3]{\log_2x}&=\frac{\log_2x}{3} \end{align*}
Now setting $y=\log_2x$, we have
 \[\sqrt[3]{y}=\frac{y}{3}\]Solving gets $y=\log_2x=3\sqrt{3}\Longrightarrow(\log_2x)^2=(3\sqrt{3})^2=\boxed{27}$.

 



\textbf{Solution 4}

Say that $\log_{2^3}x=a$ and $\log_2x=b$ so we have $\log_2a=\log_{2^3}b$. And we want $b^2$.

 $\\ \log_2a=\frac13 \log_{2}b \ \ \text{\tiny{(step 1)}}\\ \frac{\log_2a}{\log_{2}b}=\log_ba=\frac13\\ b^{1/3}=a.$

 Because $3a=b$ (as $2^{3a}=x$ and $2^b=x$ from our setup), we have that 

 $b^{1/3}=\frac{b}{3}\\ b^{-2/3}=\frac13\\ b=3^{3/2}\\ \\b^2=3^3=\boxed{27}$

 Note that we use the property $\log_{x^n}y=\frac1n\log_xy$ in step 1 and $\frac{\log_wx}{\log_wy}=\log_yx$ in step 2 in this solution.

 


	\item (\textit{AIME I 1988}) Suppose that $|x_i| < 1$ for $i = 1, 2, \dots, n$. Suppose further that


 $|x_1| + |x_2| + \dots + |x_n| = 19 + |x_1 + x_2 + \dots + x_n|$.

 What is the smallest possible value of $n$?


 
 \textbf{Answer: }020



\textbf{Solution}

Since $|x_i| < 1$ then

 
 \[|x_1| + |x_2| + \dots + |x_n| = 19 + |x_1 + x_2 + \dots + x_n| < n.\]

 So $n \ge 20$. We now just need to find an example where $n = 20$: suppose $x_{2k-1} = \frac{19}{20}$ and $x_{2k} = -\frac{19}{20}$; then on the left hand side we have $\left|\frac{19}{20}\right| + \left|-\frac{19}{20}\right| + \dots + \left|-\frac{19}{20}\right| = 20\left(\frac{19}{20}\right) = 19$. On the right hand side, we have $19 + \left|\frac{19}{20} - \frac{19}{20} + \dots - \frac{19}{20}\right| = 19 + 0 = 19$, and so the equation can hold for $n = \boxed{020}$.

 


 



\textbf{Solution 2 (Motivating solution)}

First off, one can test $1,-1,1$ and find that the $LHS$ is $3$ and the RHS is $1.$ Similarly testing $1,-1,-1,1$ yields $4$ on the LHS and $0$ on the RHS. It seems for every negative we gain twice of that negative on the LHS. However, when we test something like $1,-1,-1,-1,1$ we find the LHS to be $5$ and the RHS to be $1.$ What happened? There were more negatives than positives. Why does this mean that the LHS doesn't grow? There aren't enough positives to "cancel out!" Therefore if for every negative we need a positive for it to cancel out to grow. We can make the $LHS$ grow by approximately $2$ every time we choose a negative and a positive, so if we have $20$ numbers, namely $10$ positive and $10$ negative we can obtain the desired answer.

 



\textbf{Solution 3}

Straight off, we notice that the RHS must be greater than or equal to 19, because an absolute value only gives nonnegative values. It then becomes clear that $|x_1|+|x_2|+...+|x_n|\ge19$. If each $x_n$ were equal to 1, then $n=19$. However, $x_n<1$, so there must be at least one extra term to satisfy the inequality. Therefore, $n=\boxed{020}$.

 



\textbf{Solution 4 (Similar to 3)}

Note as long as $x_i < 1$ then we can set two terms $x_i, x_j$ as opposites on the number line and the result follows. That then implies for all even $n$ that $|x_1+\dots+x_n| + 19 = 19$ hence $\sum |x_i| \geq 19$. Equality holds if all $x_i = 1$ but we have the property $x_i < 1$. So, we must account for this with one more term $x_i$ making $n = \boxed{20} \quad \square$.

 


	\item (\textit{AIME I 1988}) Let $\frac{m}{n}$, in lowest terms, be the probability that a randomly chosen positive divisor of $10^{99}$ is an integer multiple of $10^{88}$. Find $m + n$.


 
 \textbf{Answer: }634



\textbf{Solution}

$10^{99} = 2^{99}5^{99}$, so it has $(99 + 1)(99 + 1) = 10000$ factors. Out of these, we only want those factors of $10^{99}$ which are divisible by $10^{88}$; it is easy to draw a \href{https://artofproblemsolving.com/wiki/index.php?title=Bijection}{bijection} to the number of factors that $10^{11} = 2^{11}5^{11}$ has, which is $(11 + 1)(11 + 1) = 144$. Our probability is $\frac{m}{n} = \frac{144}{10000} = \frac{9}{625}$, and $m + n = \boxed{634}$.

 



\textbf{Solution 2}

Like before there are \( (99+1)(99+1) = 100^2 = 10,000 \) positive divisors of \( 10^{99} \). Now we calculate the cases that satisfy the conditions. Notice that every multiple must be in the form \( 2^i 5^j \), where \( 88 \le i, j \le 99 \). Then, there are \( 12 \) choices for both \( i \) and \( j \). We can have any combination of \( i \) and \( j \) together, giving us \( 144 \) favorable outcomes. Thus our probability is \( \frac{144}{10,000} \), or \( \frac{9}{625} \), in which \( \frac{m}{n} = \frac{9}{625} \). Thus, \( m+n = 9+625 = \boxed{634} \)

 


	\item (\textit{AIME I 1988}) It is possible to place positive integers into the vacant twenty-one squares of the 5 times 5 square shown below so that the numbers in each row and column form arithmetic sequences. Find the number that must occupy the vacant square marked by the asterisk (*).


 
                    \begin{figure}[H]
                        \centering
                        \includegraphics[width=0.5\textwidth]{Problem/AIME_I_1988_6_Problem_Figure_1.png}
                    \end{figure}
                    

 \textbf{Answer: }142



\textbf{Solution 1 (specific)}

Let the coordinates of the square at the bottom left be $(0,0)$, the square to the right $(1,0)$, etc.

 Label the leftmost column (from bottom to top) $0, a, 2a, 3a, 4a$ and the bottom-most row (from left to right) $0, b, 2b, 3b, 4b$. Our method will be to use the given numbers to set up equations to solve for $a$ and $b$, and then calculate $(*)$.

 $\begin{tabular}[b]{|c|c|c|c|c|}\hline 4a & & & * & \\ \hline 3a & 74 & & & \\ \hline 2a & & & & 186 \\ \hline a & & 103 & & \\ \hline 0 & b & 2b & 3b & 4b \\ \hline \end{tabular}$

 We can compute the squares at the intersections of two existing numbers in terms of $a$ and $b$; two such equations will give us the values of $a$ and $b$. On the fourth row from the bottom, the common difference is $74 - 3a$, so the square at $(2,3)$ has a value of $148 - 3a$. On the third column from the left, the common difference is $103 - 2b$, so that square also has a value of $2b + 3(103 - 2b) = 309 - 4b$. Equating, we get $148 - 3a = 309 - 4b \Longrightarrow 4b - 3a = 161$.

 Now we compute the square $(2,2)$. By rows, this value is simply the average of $2a$ and $186$, so it is equal to $\frac{2a + 186}{2} = a + 93$. By columns, the common difference is $103 - 2b$, so our value is $206 - 2b$. Equating, $a + 93 = 206 - 2b \Longrightarrow a + 2b = 113$.

 Solving 
 \begin{align*}4b - 3a &= 161\\ a + 2b &= 113 \end{align*}


 gives $a = 13$, $b = 50$. Now it is simple to calculate $(4,3)$. One way to do it is to see that $(2,2)$ has $206 - 2b = 106$ and $(4,2)$ has $186$, so $(3,2)$ has $\frac{106 + 186}{2} = 146$. Now, $(3,0)$ has $3b = 150$, so $(3,2) = \frac{(3,0) + (3,4)}{2} \Longrightarrow (3,4) = * = \boxed{142}$.

 



\textbf{Solution 2 (general)}

First, let $a =$ the number to be placed in the first column, fourth row.  Let $b =$ the number to be placed in the second column, fifth row.  We can determine the entire first column and fifth row in terms of $a$ and $b$:

 $\begin{tabular}[b]{|c|c|c|c|c|}\hline 4a & & & & \\ \hline 3a & & & & \\ \hline 2a & & & & \\ \hline a & & & & \\ \hline 0 & b & 2b & 3b & 4b \\ \hline \end{tabular}$

 Next, let $a + b + c =$ the number to be placed in the second column, fourth row.  We can determine the entire second column and fourth row in terms of $a$, $b$, and $c$:

 $\begin{tabular}[b]{|c|c|c|c|c|}\hline 4a & 4a + b + 4c & & & \\ \hline 3a & 3a + b + 3c & & & \\ \hline 2a & 2a + b + 2c & & & \\ \hline a & a + b + c & a + 2b + 2c & a + 3b + 3c & a + 4b + 4c \\ \hline 0 & b & 2b & 3b & 4b \\ \hline \end{tabular}$

 We have now determined at least two values in each row and column.  We can finish the table without introducing any more variables:

 $\begin{tabular}[b]{|c|c|c|c|c|}\hline 4a & 4a + b + 4c & 4a + 2b + 8c & 4a + 3b + 12c & 4a + 4b + 16c \\ \hline 3a & 3a + b + 3c & 3a + 2b + 6c & 3a + 3b + 9c & 3a + 4b + 12c \\ \hline 2a & 2a + b + 2c & 2a + 2b + 4c & 2a + 3b + 6c & 2a + 4b + 8c \\ \hline a & a + b + c & a + 2b + 2c & a + 3b + 3c & a + 4b + 4c \\ \hline 0 & b & 2b & 3b & 4b \\ \hline \end{tabular}$

 We now have a system of equations.

 <div style="text-align:center;">$3a + b + 3c = 74$
$2a + 4b + 8c = 186$
$a + 2b + 2c = 103$</div>Solving, we find that $(a,b,c) = (13,50, - 5)$.  The number in the square marked by the asterisk is $4a + 3b + 12c = \boxed{142}$




\textbf{Solution 3 (Only one variable)}

We begin with the table that was given to us and add in the following arithmetic progression on the bottom:

 $\begin{tabular}[b]{|c|c|c|c|c|}\hline & & & * & \\ \hline & 74 & & & \\ \hline & & & & 186 \\ \hline & & 103 & & \\ \hline 0 & x & 2x & 3x & 4x \\ \hline \end{tabular}$

 Since all the rows and columns satisfy an arithmetic progression, we have the following:

 $\begin{tabular}[b]{|c|c|c|c|c|}\hline & & 412 - 6x & 392 - 5x & 372 - 4x \\ \hline & 74 & 309 - 4x  & 294 - 3x  & 279 - 2x \\ \hline & & 206 - 2x & 196 - x & 186 \\ \hline & & 103 & 98 + x & 2x + 93 \\ \hline 0 & x & 2x & 3x & 4x \\ \hline \end{tabular}$

 We can solve for $x$ in the 2nd row, namely $324 - 5x = 74$ because the arithmetic progression from left to right has difference $x - 15$. Therefore, we have $x = 50$, and because the desired asterisk is $392 - 5x$, the answer is $392 - 250$ = $\boxed {142}$.

 Note: Since $324-5x=74,$ we can just say that $392-5x = 74 + (324-5x) = 74 + 68,$ which gets us the same answer of $\boxed{142}.$ ~PrimeParadox (also thank you Arcticturn, this solution really helped)

 


	\item (\textit{AIME I 1988}) In triangle $ABC$, $\tan \angle CAB = 22/7$, and the altitude from $A$ divides $BC$ into segments of length $3$ and $17$. What is the area of triangle $ABC$?


 
 \textbf{Answer: }110



\textbf{Solution}


                    \begin{figure}[H]
                        \centering
                        \includegraphics[width=0.5\textwidth]{Solution/AIME_I_1988_7_Solution_Figure_1.png}
                    \end{figure}
                    

 Call $\angle BAD$ $\alpha$ and $\angle CAD$ $\beta$. So, $\tan \alpha = \frac {17}{h}$ and $\tan \beta = \frac {3}{h}$. Using the tangent addition formula $\tan (\alpha + \beta) = \dfrac {\tan \alpha + \tan \beta}{1 - \tan \alpha \cdot \tan \beta}$, we get $\tan (\alpha + \beta) = \dfrac {\frac {20}{h}}{\frac {h^2 - 51}{h^2}} = \frac {22}{7}$. 

 Simplifying, we get $\frac {20h}{h^2 - 51} = \frac {22}{7}$. Cross-multiplying and simplifying, we get $11h^2-70h-561 = 0$. Factoring, we get $(11h+51)(h-11) = 0$, so we take the positive positive solution, which is $h = 11$. Therefore, the answer is $\frac {20 \cdot 11}{2} = 110$, so the answer is $\boxed {110}$.

 


	\item (\textit{AIME I 1988}) The function $f$, defined on the set of ordered pairs of positive integers, satisfies the following properties:

 \[f(x, x) = x,\; f(x, y) = f(y, x), {\textrm \ and\ } (x+y)f(x, y) = yf(x, x+y).\]
Calculate $f(14,52)$.


 
 \textbf{Answer: }364



\textbf{Solution 1 (Algebra)}

Let $z = x+y$. By the substitution $z=x+y,$ we rewrite the third property in terms of $x$ and $z,$ then solve for $f(x,z):$
 \begin{align*} zf(x,z-x) &= (z-x)f(x,z) \\ f(x,z) &= \frac{z}{z-x} \cdot f(x,z-x). \end{align*}
Using the properties of $f,$ we have
 \begin{align*} f(14,52) &= \frac{52}{38} \cdot f(14,38) \\ &= \frac{52}{38} \cdot \frac{38}{24} \cdot f(14,24) \\ &= \frac{52}{38} \cdot \frac{38}{24} \cdot \frac{24}{10} \cdot f(14,10)\\ &= \frac{52}{38} \cdot \frac{38}{24} \cdot \frac{24}{10} \cdot f(10,14)\\ &= \frac{52}{38} \cdot \frac{38}{24} \cdot \frac{24}{10} \cdot \frac{14}{4} \cdot f(10,4)\\ &= \frac{52}{38} \cdot \frac{38}{24} \cdot \frac{24}{10} \cdot \frac{14}{4} \cdot f(4,10)\\ &= \frac{52}{38} \cdot \frac{38}{24} \cdot \frac{24}{10} \cdot \frac{14}{4} \cdot \frac{10}{6} \cdot f(4,6)\\ &= \frac{52}{38} \cdot \frac{38}{24} \cdot \frac{24}{10} \cdot \frac{14}{4} \cdot \frac{10}{6} \cdot \frac{6}{2} \cdot f(4,2)\\ &= \frac{52}{38} \cdot \frac{38}{24} \cdot \frac{24}{10} \cdot \frac{14}{4} \cdot \frac{10}{6} \cdot \frac{6}{2} \cdot f(2,4)\\ &= \frac{52}{38} \cdot \frac{38}{24} \cdot \frac{24}{10} \cdot \frac{14}{4} \cdot \frac{10}{6} \cdot \frac{6}{2} \cdot \frac{4}{2} \cdot f(2,2)\\ &= \frac{52}{38} \cdot \frac{38}{24} \cdot \frac{24}{10} \cdot \frac{14}{4} \cdot \frac{10}{6} \cdot \frac{6}{2} \cdot \frac{4}{2} \cdot 2\\ &=\boxed{364}. \end{align*}
~MRENTHUSIASM (credit given to AoPS)

 



\textbf{Solution 2 (Algebra)}

Since all of the function's properties contain a recursive definition except for the first one, we know that $f(x,x) = x$ in order to obtain an integer answer. So, we have to transform $f(14,52)$ to this form by exploiting the other properties. The second one doesn't help us immediately, so we will use the third one. 

 Note that 
 \[f(14,52) = f(14,14 + 38) = \frac{52}{38}\cdot f(14,38).\]

 Repeating the process several times,
 \begin{align*} f(14,52) & = f(14,14 + 38) \\ & = \frac{52}{38}\cdot f(14,38) \\ & = \frac{52}{38}\cdot \frac{38}{24}\cdot f(14,14 + 24) \\ & = \frac{52}{24}\cdot f(14,24) \\ & = \frac{52}{10}\cdot f(10,14) \\ & = \frac{52}{10}\cdot \frac{14}{4}\cdot f(10,4) \\ & = \frac{91}{5}\cdot f(4,10) \\ & = \frac{91}{3}\cdot f(4,6) \\ & = 91\cdot f(2,4) \\ & = 91\cdot 2 \cdot f(2,2) \\ & = \boxed{364}. \end{align*}





\textbf{Solution 3 (Number Theory)}

Notice that $f(x,y) = \mathrm{lcm}(x,y)$ satisfies all three properties: 

 For the first two properties, it is clear that $\mathrm{lcm}(x,x) = x$ and $\mathrm{lcm}(x,y) = \mathrm{lcm}(y,x)$. 

 For the third property, using the identities $\gcd(x,y) \cdot \mathrm{lcm}(x,y) = x\cdot y$ and $\gcd(x,x+y) = \gcd(x,y)$ gives
 \begin{align*} y \cdot \mathrm{lcm}(x,x+y) &= \dfrac{y \cdot x(x+y)}{\gcd(x,x+y)} \\ &= \dfrac{(x+y) \cdot xy}{\gcd(x,y)} \\ &= (x+y) \cdot \mathrm{lcm}(x,y). \end{align*}
Hence, $f(x,y) = \mathrm{lcm}(x,y)$ is a solution to the functional equation. 

 Since this is an AIME problem, there is exactly one correct answer, and thus, exactly one possible value of $f(14,52)$. 

 Therefore, we have 
 \begin{align*} f(14,52) &= \mathrm{lcm}(14,52) \\ &= \mathrm{lcm}(2 \cdot 7,2^2 \cdot 13) \\ &= 2^2 \cdot 7 \cdot 13 \\ &= \boxed{364}. \end{align*}


 Proof that $f(x,y)=\mathrm{lcm}(x,y)$ is the only function instead of just using the fact that this is an AIME problem (Proof credited to AMSP's 101 Algebra Problems):

 FTSOC, suppose that there is another function $g(x,y)$ also satisfying the given conditions. Let $S$ be the set of all pairs of positive integers $(x,y)$ such that $f(x,y) \ne g(x,y)$, and let $(m,n)$ be such a pair with minimal sum $m+n$. It is clear that $m \ne n$, otherwise $f(m,n) = f(m,m) = m = g(m,m) = g(m,n)$. By symmetry, we can also assume that $n>m$. Note that 
 \[nf(m,n-m)=[m+(n-m)]f(m,n-m)\]
 \[=(n-m)f(m,m+(n-m))=(n-m)f(m,n)\]or 
 \[f(m,n-m=\dfrac{n-m}{n} \cdot f(m,n).\]Likewise, 
 \[g(m,n-m) = \dfrac{n-m}{n} \cdot g(m,n).\]Since $f(m,n) \ne g(m,n)$, $f(m,n-m) \ne g(m,n-m)$. Thus $(m,n-m) \in S$. But $(m,n-m)$ has a smaller sum then $(m,n)$, a contradiction. Therefore, $f(x,y) = \mathrm{lcm}(x,y)$ is the only solution.

 


	\item (\textit{AIME I 1988}) Find the smallest positive integer whose \href{https://artofproblemsolving.com/wiki/index.php?title=Perfect\_cube}{cube} ends in $888$.


 
 \textbf{Answer: }192



\textbf{Solution 1}

A little bit of checking tells us that the units digit must be 2. Now our cube must be in the form of $(10k + 2)^3$; using the \href{https://artofproblemsolving.com/wiki/index.php?title=Binomial\_theorem}{binomial theorem} gives us $1000k^3 + 600k^2 + 120k + 8$. Since we are looking for the tens digit, $\mod{100}$ we get $20k + 8 \equiv 88 \pmod{100}$. This is true if the tens digit is either $4$ or $9$. Casework:

 
            \begin{itemize}
        \item $4$: Then our cube must be in the form of $(100k + 42)^3 \equiv 3(100k)(42)^2 + 42^3 \equiv 200k + 88 \pmod{1000}$. Hence the lowest possible value for the hundreds digit is $4$, and so $442$ is a valid solution.
        \item $9$: Then our cube is $(100k + 92)^3 \equiv 3(100k)(92)^2 + 92^3 \equiv 200k + 688 \pmod{1000}$. The lowest possible value for the hundreds digit is $1$, and we get $192$. Hence, since $192 < 442$, the answer is $\fbox{192}$
            \end{itemize}
            



\textbf{Solution 2}

$n^3 \equiv 888 \pmod{1000} \implies n^3 \equiv 0 \pmod 8$ and $n^3 \equiv 13 \pmod{125}$.$n \equiv 2 \pmod 5$ due to the last digit of $n^3$. Let $n = 5a + 2$. By expanding, $125a^3 + 150a^2 + 60a + 8 \equiv 13 \pmod{125} \implies 5a^2 + 12a \equiv 1 \pmod{25}$. 

 By looking at the last digit again, we see $a \equiv 3 \pmod5$, so we let $a = 5a_1 + 3$ where $a_1 \in \mathbb{Z^+}$. Plugging this in to $5a^2 + 12a \equiv 1 \pmod{25}$ gives $10a_1 + 6 \equiv 1 \pmod{25}$. Obviously, $a_1 \equiv 2 \pmod 5$, so we let $a_1 = 5a_2 + 2$ where $a_2$ can be any non-negative integer. 

 Therefore, $n = 2 + 5(3+ 5(2+5a_2)) = 125a_2 + 67$. $n^3$ must also be a multiple of $8$, so $n$ must be even. $125a_2 + 67 \equiv 0 \pmod 2 \implies a_2 \equiv 1 \pmod 2$. Therefore, $a_2 = 2a_3 + 1$, where $a_3$ is any non-negative integer. The number $n$ has form $125(2a_3+1)+67 = 250a_3+192$. So the minimum $n = \boxed{192}$.

 BTW, explanation is pretty bad

 



\textbf{Solution 3}

Let $x^3 = 1000a + 888$. We factor an $8$ out of the right hand side, and we note that $x$ must be of the form $x = 2y$, where $y$ is a positive integer. Then, this becomes $y^3 = 125a + 111$. Taking mod $5$, $25$, and $125$, we get $y^3 \equiv 1\pmod 5$, $y^3 \equiv 11\pmod{25}$, and $y^3 \equiv 111\pmod{125}$. 

 We can work our way up, and find that $y\equiv 1\pmod 5$, $y\equiv 21\pmod{25}$, and finally $y\equiv 96\pmod{125}$. This gives us our smallest value, $y = 96$, so $x = \boxed{192}$, as desired. - Spacesam

 



\textbf{Solution 4 (Bash)}

Let this integer be $x.$ Note that 
 \[x^3 \equiv 888 \pmod{1000} \implies x \equiv 0 \pmod {2}~~ \cap ~~ x \equiv 2 \pmod{5}.\] We wish to find the residue of $x$ mod $125.$ Note that 
 \[x \equiv 2,7,12,17, \text{ or } 22 \pmod{25}\] using our congruence in mod $5.$ The residue that works must also satisfy $x^3 \equiv 13 \pmod{25}$ from our original congruence. Noting that $17^3 \equiv (-8)^3 \equiv -512 \equiv 13 \pmod{25}$ (and bashing out the other residues perhaps but they're not that hard), we find that 
 \[x \equiv 17 \pmod{25}.\] Observe that 
 \[17^3 \equiv \mathbf{4913} \equiv 38 \pmod{125}.\] Also, 
 \[42^3 \equiv 42^2 \cdot 42 \equiv 1764 \cdot 42 \equiv 14 \cdot 42 \equiv 88 \pmod{125}.\] Finally, 
 \[67^3 \equiv 67^2 \cdot 67 \equiv 4489 \cdot 67 \equiv (-11) \cdot 67 \equiv -737 \equiv 13 \pmod{125},\] as desired. Thus, $x$ must satisfy 
 \[x \equiv 0 \pmod{2}~~ \cap ~~x \equiv 67 \pmod{125} \implies x \equiv 192 \pmod{250} \implies x=\boxed{192}.\] ~samrocksnature

 



\textbf{Solution 5}

This number is in the form of $10k+2$, after binomial expansion, we only want $600k^2+120k\equiv 880 \pmod{1000}$. We realize that $600,120$ are both multiples of $8$, we only need that $600k^2+120k \equiv 5\pmod{125}$, so we write $600k^2+120k=125x+5; 120k^2+24k=25x+1, 24(5k^2+k)=25x+1, 5k^2+k\equiv -1\pmod{25}$

 Then, we write $5k^2+k=25m-1, 5k^2+k+1=25m$ so $k+1$ must be a multiple of $5$ at least, so $k\equiv {-1, -6, -11, -16, -21} \pmod {25}$ after checking, when $k=-6, 5k^2+k+1=175=25\cdot 7$. So $k\equiv -6 \pmod{25}$, smallest $k=19$, the number is $\boxed{192}$




\textbf{Solution 6 (A bit of brute force using basic knowledge.)}

We do know the unit digit has to be 2, So lets consider the number of the form of $x2$.

 On cubing $x2$, we get a number of the form $x^3 \,6x^2 \,12x \,8$ where the unit digit of $12x$ must be 8, therefore x can be 4 or 9. But for this value of $x$ the hundreds digit won't be 8.

 Thus our number must be of the form $x42$ / $x92$. Repeating the above process we get values of $x$ as 4 and 1 respectively.

 Therefore the smallest number is 192.~ Dubey619

 


	\item (\textit{AIME I 1988}) A convex polyhedron has for its faces 12 squares, 8 regular hexagons, and 6 regular octagons. At each vertex of the polyhedron one square, one hexagon, and one octagon meet. How many segments joining vertices of the polyhedron lie in the interior of the polyhedron rather than along an edge or a face?


 
 \textbf{Answer: }840



\textbf{Solution 1}

The polyhedron described looks like this, a truncated cuboctahedron.

 
                    \begin{figure}[H]
                        \centering
                        \includegraphics[width=0.5\textwidth]{Solution/AIME_I_1988_10_Solution_Figure_1.jpg}
                    \end{figure}
                    The number of segments joining the vertices of the polyhedron is ${48\choose2} = 1128$. We must now subtract out those segments that lie along an edge or a face.

 Since every vertex of the polyhedron lies on exactly one vertex of a square/hexagon/octagon, we have that $V = 12 \cdot 4 = 8 \cdot 6 = 6 \cdot 8 = 48$.

 Each vertex is formed by the intersection of 3 edges. Since every edge is counted twice, once at each of its endpoints, the number of edges $E$ is $\frac{3}{2}V = 72$.

 Each of the segments lying on a face of the polyhedron must be a diagonal of that face. Each square contributes $\frac{n(n-3)}{2} = 2$ diagonals, each hexagon $9$, and each octagon $20$. The number of diagonals is thus $2 \cdot 12 + 9 \cdot 8 + 20 \cdot 6 = 216$.

 Subtracting, we get that the number of space diagonals is $1128 - 72 - 216 = \boxed{840}$.

 



\textbf{Solution 2}

We first find the number of vertices on the polyhedron:There are 4 corners per square, 6 corners per hexagon, and 8 corners per octagon. Each vertex is where 3 corners coincide, so

 


	\item (\textit{AIME I 1988}) Let $w_1, w_2, \dots, w_n$ be complex numbers. A line $L$ in the complex plane is called a mean line for the points $w_1, w_2, \dots, w_n$ if $L$ contains points (complex numbers) $z_1, z_2, \dots, z_n$ such that


 $\sum_{k = 1}^n (z_k - w_k) = 0.$

 For the numbers $w_1 = 32 + 170i$, $w_2 = -7 + 64i$, $w_3 = -9 +200i$, $w_4 = 1 + 27i$, and $w_5 = -14 + 43i$, there is a unique mean line with y-intercept $3$. Find the slope of this mean line.


 
 \textbf{Answer: }163



\textbf{Solution 1}

$\sum_{k=1}^5 z_k - \sum_{k=1}^5 w_k = 0$
$\sum_{k=1}^5 z_k = 3 + 504i$

 Each $z_k = x_k + y_ki$ lies on the complex line $y = mx + 3$, so we can rewrite this as 

 $\sum_{k=1}^5 z_k = \sum_{k=1}^5 x_k + \sum_{k=1}^5 y_ki$
$3 + 504i = \sum_{k=1}^5 x_k + i \sum_{k=1}^5 (mx_k + 3)$

 Matching the real parts and the imaginary parts, we get that $\sum_{k=1}^5 x_k = 3$ and $\sum_{k=1}^5 (mx_k + 3) = 504$. Simplifying the second summation, we find that $m\sum_{k=1}^5 x_k = 504 - 3 \cdot 5 = 489$, and substituting, the answer is $m \cdot 3 = 489 \Longrightarrow m = 163$.

 



\textbf{Solution 2}

We know that 

 $\sum_{k=1}^5 w_k = 3 + 504i$

 And because the sum of the 5 $z$'s must cancel this out,

 $\sum_{k=1}^5 z_k = 3 + 504i$

 We write the numbers in the form $a + bi$ and we know that 

 $\sum_{k=1}^5 a_k = 3$ and $\sum_{k=1}^5 b_k = 504$

 The line is of equation $y=mx+3$. Substituting in the coordinates, we have $b_k = ma_k + 3$.

 Summing all 5 of the equations given for each $k$, we get 

 $504 =  3m + 15$

 Solving for $m$, the slope, we get $\boxed{163}$




\textbf{Solution 3}

The mean line for $w_1, . . ., w_5$ must pass through the mean (the center of mass) of these points, which, if we graph them on the complex plane, is $(\frac{3}{5}, \frac{504i}{5})$. Since we now have two points, namely that one and $(0, 3i)$, we can simply find the slope between them, which is $\boxed{163}$ by the good ol' slope formula.

 


	\item (\textit{AIME I 1988}) Let $P$ be an interior point of triangle $ABC$ and extend lines from the vertices through $P$ to the opposite sides. Let $a$, $b$, $c$, and $d$ denote the lengths of the segments indicated in the figure. Find the product $abc$ if $a + b + c = 43$ and $d = 3$. 


 
                    \begin{figure}[H]
                        \centering
                        \includegraphics[width=0.5\textwidth]{Problem/AIME_I_1988_12_Problem_Figure_1.png}
                    \end{figure}
                    

 \textbf{Answer: }441



\textbf{Solution 1 (Ceva's Theorem sum form)}

Call the \href{https://artofproblemsolving.com/wiki/index.php?title=Cevian}{cevians} AD, BE, and CF. Using area ratios ($\triangle PBC$ and $\triangle ABC$ have the same base), we have: 

 $\frac {d}{a + d} = \frac {[PBC]}{[ABC]}$

 Similarily, $\frac {d}{b + d} = \frac {[PCA]}{[ABC]}$ and $\frac {d}{c + d} = \frac {[PAB]}{[ABC]}$.

 Then,$\frac {d}{a + d} + \frac {d}{b + d} + \frac {d}{c + d} = \frac {[PBC]}{[ABC]} + \frac {[PCA]}{[ABC]} + \frac {[PAB]}{[ABC]} = \frac {[ABC]}{[ABC]} = 1$

 The \href{https://artofproblemsolving.com/wiki/index.php?title=Identity}{identity} $\frac {d}{a + d} + \frac {d}{b + d} + \frac {d}{c + d} = 1$ is a form of \href{https://artofproblemsolving.com/wiki/index.php?title=Ceva\%27s\_Theorem}{Ceva's Theorem}. 

 Plugging in $d = 3$, we get 

 
 \[\frac{3}{a + 3} + \frac{3}{b + 3} + \frac{3}{c+3} = 1\]
 \[3[(a + 3)(b + 3) + (b + 3)(c + 3) + (c + 3)(a + 3)] = (a+3)(b+3)(c+3)\]
 \[3(ab + bc + ca) + 18(a + b + c) + 81 = abc + 3(ab + bc + ca) + 9(a + b + c) + 27\]
 \[9(a + b + c) + 54 = abc=\boxed{441}\]



	\item (\textit{AIME I 1988}) Find $a$ if $a$ and $b$ are integers such that $x^2 - x - 1$ is a factor of $ax^{17} + bx^{16} + 1$.


 
 \textbf{Answer: }987



\textbf{Solution 1 (Fibonacci Numbers)}

Let $F_n$ represent the $n$th number in the Fibonacci sequence. Therefore,
 \begin{align*} x^2 - x - 1 = 0&\Longrightarrow x^n = F_n(x), \ n\in N \\ &\Longrightarrow x^{n + 2} = F_{n + 1}\cdot x + F_n,\ n\in N. \end{align*}
The above uses the similarity between the Fibonacci recursion|recursive definition, $F_{n+2} - F_{n+1} - F_n = 0$, and the polynomial $x^2 - x - 1 = 0$. 
 \begin{align*} 0 = ax^{17} + bx^{16} + 1 = a(F_{17}\cdot x + F_{16}) + b(F_{16}\cdot x + F_{15}) + 1 &\Longrightarrow (aF_{17} + bF_{16})\cdot x + (aF_{16} + bF_{15} + 1) = 0,\ x\not\in Q \\ &\Longrightarrow aF_{17} + bF_{16} = 0 \text{ and } aF_{16} + bF_{15} + 1 = 0 \\ &\Longrightarrow a = F_{16},\ b = - F_{17} \\ &\Longrightarrow a = \boxed {987}. \end{align*}


 \textbf{Note} Feel free to correct if I missed something, but its like late my time sooo...

 Actually, the author accidently shifted the indices by 1! Indeed, $a = F_{15}$ here by the definition, not $F_{16}$. The only reason this still works is because we are under the assumption $F_0 = 1, F_1 = 1$, as $x^2 = x+1$. 

 For those not following, we have that $x^{n+2} = F_{n+1}x + F_n$. Hence, $ax^{17} + bx^{16} + 1 = a(F_{16}x + F_{15}) + b(F_{15}x + F_{14}) + 1 \implies a = F_{15}$.

 



\textbf{Solution 2 (Fibonacci Numbers)}

We can long divide and search for a pattern; then the remainder would be set to zero to solve for $a$. Writing out a few examples quickly shows us that the remainders after each subtraction follow the Fibonacci sequence. Carrying out this pattern, we find that the remainder is 
 \[(F_{16}b + F_{17}a)x + F_{15}b + F_{16}a + 1 = 0.\] Since the coefficient of $x$ must be zero, this gives us two equations, $F_{16}b + F_{17}a = 0$ and $F_{15}b + F_{16}a + 1 = 0$. Solving these two as above, we get that $a = \boxed{987}$.

 There are various similar solutions which yield the same pattern, such as repeated substitution of $x^2 = x + 1$ into the polynomial with a higher degree, as shown in Solution 6.

 



\textbf{Solution 3 (Fibonacci Numbers: For Beginners, Less Technical)}

Trying to divide $ax^{17} + bx^{16} + 1$ by $x^2-x-1$ would be very tough, so let's try to divide using smaller degrees of $x$. Doing $\frac{ax^3+bx^2+1}{x^2-x-1}$, we get the following systems of equations:
 \begin{align*} a+b &= -1, \\ 2a+b &= 0. \end{align*}
Continuing with $\frac{ax^4+bx^3+1}{x^2-x-1}$:
 \begin{align*} 2a+b &= -1, \\ 3a+2b &= 0. \end{align*}
There is somewhat of a pattern showing up, so let's try $\frac{ax^5+bx^4+1}{x^2-x-1}$ to verify. We get:
 \begin{align*} 3a+2b &= -1, \\ 5a+3b &= 0. \end{align*}
Now we begin to see that our pattern is actually the Fibonacci Numbers! Using the previous equations, we can make a general statement about $\frac{ax^n+bx^{n-1}+1}{x^2-x-1}$:
 \begin{align*} af_{n-1}+bf_{n-2} &= -1, \\ af_n+bf_{n-1} &= 0. \end{align*}
Also, noticing our solutions from the previous systems of equations, we can create the following statement:

 If $ax^n+bx^{n-1}+1$ has $x^2-x-1$ as a factor, then $a=f_{n-1}$ and $b = f_n.$

 Thus, if $ax^{17}+bx^{16}+1$ has $x^2-x-1$ as a factor, we get that $a = 987$ and $b = -1597,$ so $a = \boxed {987}$.

 



\textbf{Solution 4 (Fibonacci Numbers: Not Rigorous)}

Let's work backwards! Let $F(x) = ax^{17} + bx^{16} + 1$ and let $P(x)$ be the polynomial such that $P(x)(x^2 - x - 1) = F(x)$.

 Clearly, the constant term of $P(x)$ must be $- 1$. Now, we have 
 \[(x^2 - x - 1)(c_1x^{15} + c_2x^{14} + \cdots + c_{15}x - 1),\] where $c_{i}$ is some coefficient. However, since $F(x)$ has no $x$ term, it must be true that $c_{15} = 1$. 

 Let's find $c_{14}$ now. Notice that all we care about in finding $c_{14}$ is that $(x^2 - x - 1)(\cdots + c_{14}x^2 + x - 1) = \text{something} + 0x^2 + \text{something}$. Therefore, $c_{14} = - 2$. Undergoing a similar process, $c_{13} = 3$, $c_{12} = - 5$, $c_{11} = 8$, and we see a nice pattern. The coefficients of $P(x)$ are just the Fibonacci sequence with alternating signs! Therefore, $a = c_1 = F_{16}$, where $F_{16}$ denotes the 16th Fibonnaci number and $a = \boxed{987}$.

 



\textbf{Solution 5 (Fibonacci Numbers)}

The roots of $x^2-x-1$ are $\phi$ (the Golden Ratio) and $1-\phi$. These two must also be roots of $ax^{17}+bx^{16}+1$. Thus, we have two equations: 
 \begin{align*} a\phi^{17}+b\phi^{16}+1=0, \\ a(1-\phi)^{17}+b(1-\phi)^{16}+1=0. \end{align*}
Subtract these two and divide by $\sqrt{5}$ to get 
 \[\frac{a(\phi^{17}-(1-\phi)^{17})}{\sqrt{5}}+\frac{b(\phi^{16}-(1-\phi)^{16})}{\sqrt{5}}=0.\] Noting that the formula for the $n$th Fibonacci number is $\frac{\phi^n-(1-\phi)^n}{\sqrt{5}}$, we have $1597a+987b=0$. Since $1597$ and $987$ are coprime, the solutions to this equation under the integers are of the form $a=987k$ and $b=-1597k$, of which the only integral solutions for $a$ on $[0,999]$ are $0$ and $987$. $(a,b)=(0,0)$ cannot work since $x^2-x-1$ does not divide $1$, so the answer must be $\boxed{987}$. (Note that this solution would not be valid on an Olympiad or any test that does not restrict answers as integers between $000$ and $999$).

 



\textbf{Solution 6 (Reduces the Powers)}

We are given that $x^2 - x - 1$ is a factor of $ax^{17} + bx^{16} + 1,$ so the roots of $x^2 - x - 1$ must also be roots of $ax^{17} + bx^{16} + 1.$

 Let $x=r$ be a root of $x^2 - x - 1$ so that $r^2 - r - 1 = 0,$ or $r^2 = r + 1.$ It follows that 
 \[ar^{17} + br^{16} + 1 = 0. \hspace{20mm} (\bigstar)\]Note that
 \begin{align*} r^4 &= (r+1)^2 \\ &= r^2 + 2r + 1 \\ &= (r+1) + 2r + 1 \\ &= 3r + 2, \\ r^8 &= (3r+2)^2 \\ &= 9r^2 + 12r + 4 \\ &= 9(r+1) + 12r + 4 \\ &= 21r + 13, \\ r^{16} &= (21r + 13)^2 \\ &= 441r^2 + 546r + 169 \\ &= 441(r+1) +546r + 169 \\ &= 987r + 610. \end{align*}
We rewrite the left side of $(\bigstar)$ as a linear expression of $r:$
 \begin{align*} (ar+b)r^{16} + 1 &= 0 \\ (ar+b)(987r + 610) + 1 &= 0 \\ 987ar^2 + (610a+987b)r + 610b + 1 &= 0 \\ 987a(r+1) + (610a+987b)r + 610b + 1 &= 0 \\ (1597a+987b)r + (987a + 610b + 1) &= 0.  \end{align*}
Since this linear equation has two solutions of $r,$ it must be an identity. Therefore, we have the following system of equations:
 \begin{align*} 1597a+987b &= 0, \\ 987a+610b &= -1. \end{align*}
To eliminate $b,$ we multiply the first equation by $610$ and multiply the second equation by $987,$ then subtract the resulting equations:
 \begin{align*} 610(1597a)+610(987b) &= 0, \\ 987(987a)+987(610b) &= -987, \end{align*}
from which
 \begin{align*} (610\cdot1597-987\cdot987)a&=987 \\ (974170-974169)a&=987 \\ a&=\boxed{987}. \end{align*}
~MRENTHUSIASM

 



\textbf{Solution 7 (Uses the Roots)}

For simplicity, let $f(x) =ax^{17} + bx^{16} + 1$ and $g(x) = x^2-x-1$. Notice that the roots of $g(x)$ are also roots of $f(x)$. Let these roots be $u,v$. We get the system 
 \begin{align*} au^{17} + bu^{16} + 1 &= 0, \\ av^{17} + bv^{16} + 1 &= 0. \end{align*}
If we multiply the first equation by $v^{16}$ and the second by $u^{16}$ we get 
 \begin{align*} u^{17} v^{16} a + u^{16} v^{16} b + v^{16} &= 0, \\ u^{16} v^{17} a + u^{16} v^{16} b + u^{16} &= 0. \end{align*}
Now subtracting, we get 
 \[a(u^{17}v^{16} -u^{16} v^{17}) = u^{16}-v^{16} \implies a = \frac{u^{16} - v^{16}}{u^{17}v^{16} -u^{16} v^{17}}.\]By Vieta's, $uv=-1$ so the denominator becomes $u-v$. By difference of squares and dividing out $u-v$ we get 
 \[a= (u^8+v^8)(u^4+v^4)(u^2+v^2)(u+v).\] A simple exercise of Vieta's gets us $a= \boxed{987}.$




\textbf{Solution 8 (Engineer's Induction, only use if you don't have time left)}

We see that $ax^{17} + bx^{16} + 1 = (x^2-x-1)P(x)$ for some polynomial $P(x)$. Working forwards, we notice that the constant term of $P(x)$ must equal $-1$, to multiply the constant term of $(x^2-x-1)$ into $1$. We can similarly continue to use this logic, by repeatedly cancelling out the middle term, and obtain the process:$(x^2-x-1)(-1) = -x^2 + x + 1$$(x^2-x-1)(-1 + x) = x^3 - 2x^2 + 1$$(x^2-x-1)(-1 + x - 2x^2) = -2x^4 + 3x^3 + 1$$(x^2-x-1)(-1 + x - 2x^2 + 3x^3) = 3x^5 - 5x^4 + 1$.By this time, you can hopefully notice that the coefficient of the $x^n$ term in $P(x)$ is equal to $(-1)^{n-1} * F_{n}$, where $F_n$ equals the $n$th number in the Fibonacci sequence. From here, we just need to find the coefficient of the $x^{15}$ term in $P(x)$, which happens to be $F_{15} = \boxed{987}$.Again, try to only use Engineer's Induction when you have no other options. A rigorous proof is usually not needed, but when you have extra time, checking a solution with a rigorous method is better than worrying about your Engineer's Induction solution.

 


	\item (\textit{AIME I 1988}) Let $C$ be the graph of $xy = 1$, and denote by $C^*$ the reflection of $C$ in the line $y = 2x$. Let the equation of $C^*$ be written in the form

 \[12x^2 + bxy + cy^2 + d = 0.\]
Find the product $bc$.


 
 \textbf{Answer: }084



\textbf{Solution 1}

Given a point $P (x,y)$ on $C$, we look to find a formula for $P' (x', y')$ on $C^*$. Both points lie on a line that is \href{https://artofproblemsolving.com/wiki/index.php?title=Perpendicular}{perpendicular} to $y=2x$, so the slope of $\overline{PP'}$ is $\frac{-1}{2}$. Thus $\frac{y' - y}{x' - x} = \frac{-1}{2} \Longrightarrow x' + 2y' = x + 2y$. Also, the midpoint of $\overline{PP'}$, $\left(\frac{x + x'}{2}, \frac{y + y'}{2}\right)$, lies on the line $y = 2x$. Therefore $\frac{y + y'}{2} = x + x' \Longrightarrow 2x' - y' = y - 2x$.

 Solving these two equations, we find $x = \frac{-3x' + 4y'}{5}$ and $y = \frac{4x' + 3y'}{5}$. Substituting these points into the equation of $C$, we get $\frac{(-3x'+4y')(4x'+3y')}{25}=1$, which when expanded becomes $12x'^2-7x'y'-12y'^2+25=0$.

 Thus, $bc=(-7)(-12)=\boxed{084}$.

 



\textbf{Solution 2}

The \href{https://artofproblemsolving.com/wiki/index.php?title=Asymptote}{asymptotes} of $C$ are given by $x=0$ and $y=0$. Now if we represent the line $y=2x$ by the complex number $1+2i$, then we find the direction of the reflection of the asymptote $x=0$ by multiplying this by $2-i$, getting $4+3i$. Therefore, the asymptotes of $C^*$ are given by $4y-3x=0$ and $3y+4x=0$.

 Now to find the equation of the hyperbola, we multiply the two expressions together to get one side of the equation: $(3x-4y)(4x+3y)=12x^2-7xy-12y^2$. At this point, the right hand side of the equation will be determined by plugging the point $(\frac{\sqrt{2}}{2},\sqrt{2})$, which is unchanged by the reflection, into the expression. But this is not necessary. We see that $b=-7$, $c=-12$, so $bc=\boxed{084}$.

 


 



\textbf{Solution 3}

The matrix for a reflection about the polar line $\theta = \alpha, \alpha+\pi$ is:
 \[ \left[ \begin{array}{ccc} \cos(2\alpha) & \sin(2\alpha) \\ \sin(2\alpha) & -\cos(2\alpha) \end{array} \right] \]This is not hard to derive using a basic knowledge of linear transformations. You can refer here for more information: \href{https://en.wikipedia.org/wiki/Orthogonal\_matrix}{https://en.wikipedia.org/wiki/Orthogonal\_matrix}

 Let $\alpha = \arctan 2$. Note that the line of reflection, $y = 2x$, is the polar line $\theta = \alpha, \alpha+\pi$. Then $2\alpha = \arctan\left(-\frac{4}{3}\right)$, so $\cos(2\alpha) = -\frac{3}{5}$ and $\sin(2\alpha) = \frac{4}{5}$.

 Therefore, if $(x', y')$ is mapped to $(x, y)$ under the reflection, then $x = -\frac{3}{5}x'+\frac{4}{5}y'$ and $y = \frac{4}{5}x'+\frac{3}{5}y'$. Since the transformation matrix represents a reflection, it must be its own inverse; therefore, $x' = -\frac{3}{5}x+\frac{4}{5}y$ and $y' = \frac{4}{5}x+\frac{3}{5}y$.

 The original coordinates $(x', y')$ must satisfy $x'y' = 1$. Therefore,
 \[\left(-\frac{3}{5}x+\frac{4}{5}y\right)\left(\frac{4}{5}x+\frac{3}{5}y\right) = 1\]
 \[-\frac{12}{25}x^2 + \frac{7}{25}xy + \frac{12}{25}y^2 - 1 = 0\]
 \[12x^2 - 7xy - 12y^2 + 25 = 0\]Thus, $b = -7$ and $c = -12$, so $bc = 84$. The answer is $\boxed{084}$.

 



\textbf{Solution 4}

Find some simple points on the graph of C, reflect them and note down the new coordinates, and plug 'em into the given equation. After some plugging and chugging and solving the system of equations that follow, we get that  $b = -7$ and $c = -12$, so $bc = 84$. The answer is $\boxed{084}$.

 pi\_is\_3.141

 


	\item (\textit{AIME I 1988}) In an office at various times during the day, the boss gives the secretary a letter to type, each time putting the letter on top of the pile in the secretary's in-box. When there is time, the secretary takes the top letter off the pile and types it. There are nine letters to be typed during the day, and the boss delivers them in the order 1, 2, 3, 4, 5, 6, 7, 8, 9.


 While leaving for lunch, the secretary tells a colleague that letter 8 has already been typed, but says nothing else about the morning's typing. The colleague wonders which of the nine letters remain to be typed after lunch and in what order they will be typed. Based upon the above information, how many such after-lunch typing orders are possible? (That there are no letters left to be typed is one of the possibilities.)


 
 \textbf{Answer: }704



\textbf{Solution 1}

Re-stating the problem for clarity, let $S$ be a \href{https://artofproblemsolving.com/wiki/index.php?title=Set}{set} arranged in increasing order. At any time an element can be appended to the end of $S$, or the last element of $S$ can be removed. The question asks for the number of different orders in which all of the remaining elements of $S$ can be removed, given that $8$ had been removed already.  

 Since $8$ had already been added to the pile, the numbers $1 \ldots 7$ had already been added at some time to the pile; $9$ might or might not have been added yet. So currently $S$ is a subset of $\{1, 2, \ldots 7\}$, possibly with $9$ at the end. Given that $S$ has $k$ elements, there are $k+1$ intervals for $9$ to be inserted, or $9$ might have already been placed, giving $k+2$ different possibilities.

 Thus, the answer is $\sum_{k=0}^{7} {7 \choose k}(k+2)$ $= 1 \cdot 2 + 7 \cdot 3 + 21 \cdot 4 + 35 \cdot 5 + 35 \cdot 6 + 21 \cdot 7 + 7 \cdot 8 + 1 \cdot 9$ $= \boxed{704}$.

 \textbf{A way to compute this quickly (by AlexLikeMath)} $1 \cdot 2 + 7 \cdot 3 + 21 \cdot 4 + 35 \cdot 5 + 35 \cdot 6 + 21 \cdot 7 + 7 \cdot 8 + 1 \cdot 9$ $=1 \cdot (2+9) + 7 \cdot (3+8) + 21 \cdot (4+7) + 35 \cdot (5+6) = 64 \cdot 11 = \boxed{704}$




\textbf{Solution 2}

(MAA official solution)

 At any given time, the letters in the box are in decreasing order from top to bottom. Thus the sequence of letters in the box is uniquely determined by the set of letters in the box. We have two cases: letter 9 arrived before lunch or it did not.

 
$\textbf{Case 1:}$ Since letter 9 arrived before lunch, no further letters will arrive, and the number of possible orders is simply the number of subsets of $\mathrm{T}=\{1, 2, \ldots, 6, 7, 9\}$ which might still be in the box. In fact, each subset of $\mathrm{T}$ is possible, because the secretary might have typed letters not in the subset as soon as they arrived and not typed any others. Since $\mathrm{T}$ has 8 elements, it has $2^{8}=256$ subsets (including the empty set).

 $\textbf{Case 2:}$ Since letter 9 didn't arrive before lunch, the question is: Where can it be inserted in the typing order? Any position is possible for each subset of $\mathrm{U}=\{1,2, \ldots, 6,7\}$ which might have been left in the box during lunch (in descending order). For instance, if the letters in the box during lunch are $6,3,2,$ then the typing order 6,3,9,2 would occur if the boss would deliver letter 9 just after letter 3 was typed. There would seem to be $k+1$ places at which letter 9 could be inserted into a sequence of $\mathrm{k}$ letters. However, if letter 9 is inserted at the beginning of the sequence (i.e., at the top of the pile, so it arrives before any after-lunch typing is done), then we are duplicating an ordering from $\textbf{Case 1}$. Thus if $k$ letters are in the basket after returning from lunch, then there are $k$ places to insert letter 9 (without duplicating any $\textbf{Case 1}$ orderings). Thus we obtain
 \[\sum_{k=0}^{7} k\binom 7k=7\left(2^{7-1}\right)=448\]new orderings in $\textbf{Case 2}$.

 
Combining these cases gives $256+448=704$ possible typing orders.

 
Note. The reasoning in $\textbf{Case 2}$ can be extended to cover both cases by observing that in any sequence of $k$ letters not including letter 9, there are $k+2$ places to insert letter 9, counting the possibility of not having to insert it (i.e., if it arrived before lunch) as one of the cases. This yields
 \[\sum_{k=0}^{7}(k+2)\binom7k= \boxed{704}\]possible orderings, in agreement with the answer, found previously.

 



\textbf{Solution 3}

We can think of this as a problem of arrangements.

 Notice how, after the morning session, any of the first 7 letters could have been typed up. However, the letters that still remain must be in order from least to greatest.

 $\textbf{Case 1:}$ The 9th letter was typed up before lunch. In this case, you have $2^{7}$ possibilities for which letters have been typed and which have not.

 $\textbf{Case 2:}$ The 9th letter is to be typed up after lunch. In this case, we can take cases on the number typed up immediately before the 9th letter (still after lunch).

 Case 2.1: The 9th letter is typed first. The order of typing is now $\{9, 7, 6, 5, 4, 3, 2, 1\}$, in which any of the letters $\{1, 2\ldots7\}$ could possibly have been typed before lunch. This yields $2^7$ possibilities.

 Case 2.2: The 7th letter comes before the 9th letter. In this case, we are forcing the 7th letter to be typed after lunch. The order of typing is $\{7, 9, 6, 5, 4\ldots1\}$ where any of $\{1\ldots6\}$ could have been typed before lunch. This yields $2^6$ possibilities.

 Case 2.3: The 6th letter is typed before the 9th letter. In this case, the 6th letter is forced to be typed after lunch, however the 7th letter is not. Any of $\{7, 5, 4, 3, 2, 1\}$ could have been typed before lunch. This yields $2^6$ possibilities.

 Case 2.4: The 5th letter is typed before the 9th. The 5th and 9th letters are forced to be typed after lunch, but the others are not. For the letters $\{7, 6, 4, 3, 2, 1\}$ there are $2^6$ possibilities.

 Case 2.5: The 4th letter is typed before the 9th. There are $2^6$ possibilities for $\{7, 6, 5, 3, 2, 1\}$.

 Case 2.6: The 3rd letter is typed before the 9th. Again, there are $2^6$ possibilities.

 Case 2.7: The 2nd letter is typed before the 9th. Similarly, there are $2^6$ possibilities.

 Case 2.8: The 1st letter is typed before the 9th. There are $2^6$ possibilities.

 $\textbf{Sum:}$ The total sum is $2^7 + 2^7 + (7)2^6 = \boxed{704}$

 Clarification: For subcases of 2, the reason the letter immediately before the 9 is forced to be typed is to ensure that it is mutually exclusive from the other subcases considered

 -NiranjanR

 \textbf{Note} Problem 4.8 from the AoPS Intermediate Counting and Probability book is identical to this problem, except it is reworded and there are ten items to be dealt with (it is known that before lunch, the $9$th item has been dealt with).

 


\end{enumerate}